\SourceLeaf{244}%
given surface whose nature is expressed by the equation \(P\mathrm{d}x+Q\mathrm{d}y+R\mathrm{d}z=0\), the variations \(\delta x,\delta y,\delta z\) must also satisfy the equation \(P\delta x+Q\delta y+R\delta z=0\). It then follows readily from the usual principles that the differentials
\[
\mathrm{d}\frac{\mathrm{d}x}{\BareRoot{(\mathrm{d}x^2+\mathrm{d}y^2+\mathrm{d}z^2)}},\quad
\mathrm{d}\frac{\mathrm{d}y}{\BareRoot{(\mathrm{d}x^2+\mathrm{d}y^2+\mathrm{d}z^2)}},\quad
\mathrm{d}\frac{\mathrm{d}z}{\BareRoot{(\mathrm{d}x^2+\mathrm{d}y^2+\mathrm{d}z^2)}}
\]
must be proportional to the quantities \(P,Q,R\), respectively. Now let \(\mathrm{d}r\) be the element of the curved line, \(\lambda\) the point on the sphere representing the direction of this element, and \(L\) the point on the sphere representing the direction of the normal to the curved surface. Finally, let \(\xi,\eta,\zeta\) be the coordinates of \(\lambda\) and \(X,Y,Z\) those of \(L\), both with respect to the center of the sphere. Then we have
\[
\mathrm{d}x=\xi\mathrm{d}r,\quad\mathrm{d}y=\eta\mathrm{d}r,\quad\mathrm{d}z=\zeta\mathrm{d}r
\]
from which we conclude that those differentials become \(\mathrm{d}\xi,\mathrm{d}\eta,\mathrm{d}\zeta\). Since the quantities \(P\mathord{.}\ Q\mathord{.}\ R\) are proportional to \(X,Y,Z\), the shortest line is characterized by the equations
\[
\frac{\mathrm{d}\xi}{X}=\frac{\mathrm{d}\eta}{Y}=\frac{\mathrm{d}\zeta}{Z}
\]
It is also easy to see that \(\BareRoot{(\mathrm{d}\xi^2+\mathrm{d}\eta^2+\mathrm{d}\zeta^2)}\) equals the small arc on the sphere that measures the angle between the directions of the tangents at the beginning and end of the element \(\mathrm{d}r\). It is therefore \(=\frac{\mathrm{d}r}{\rho}\) if \(\rho\) denotes the radius of curvature of the shortest line at this point. Thus we have
\[
\rho\mathrm{d}\xi=X\mathrm{d}r,\quad\rho\mathrm{d}\eta=Y\mathrm{d}r,\quad\rho\mathrm{d}\zeta=Z\mathrm{d}r
\]
\Article{15}
Suppose that countless shortest curves issue from a given point \(A\) on a curved surface. We distinguish them by the angle that the initial element of each makes with the initial element of one of these lines chosen as the first. Let \(\varphi\) be this angle, or more generally a function of it, and let \(r\) be the length of such a shortest line from the point \(A\) to the point with coordinates \(x,y,z\). Since definite values of the variables \(r,\varphi\) correspond to definite points of the surface, the coordinates \(x,y,z\) can be regarded as functions of \(r,\varphi\). We shall retain the notation \(\lambda,L,\xi,\eta,\zeta,X,Y,Z\) in the same sense
